\subsection{Results and completion of the proof}

The complete procedure was run for every pattern in
\eqref{finite:domain}. Two implementations, described in
Appendix~\ref{app:implementation}, give the same results in
Table~\ref{finite:results}. Output $1$ denotes a contradiction from the
initial packing test; output $2$ denotes a contradiction after adding
required-path constraints or strengthening the matrix. No pattern has
output $0$.

\begin{table}[htbp]
\centering
\begin{tabular}{rrrr}
\toprule
$k$ & output $0$ & output $1$ & output $2$\\
\midrule
5&0&4&3\\
6&0&8&6\\
7&0&21&5\\
8&0&30&4\\
9&0&57&6\\
10&0&71&4\\
11&0&134&8\\
12&0&161&3\\
13&0&305&6\\
14&0&315&0\\
15&0&597&4\\
16&0&468&0\\
17&0&878&1\\
\midrule
Total&0&3049&50\\
\bottomrule
\end{tabular}
\caption{Complete finite calculation. Every pattern gives a proved
contradiction. The two successful output classes sum to $3099$.}
\label{finite:results}
\end{table}

\begin{proof}[Proof of Proposition~\ref{finite:verified}]
Suppose such a graph exists and choose an optimal path system on a
maximum clique. Lemma~\ref{finite:coverage} places its pattern in the
enumerated list. Lemma~\ref{finite:labels} gives its required-edge
representation. All assumptions used by the procedure hold: the
expansion and clique bounds are hypotheses of the proposition, the
size test follows from Lemma~\ref{size:test}, and optimality gives
Lemma~\ref{finite:extension-rule}. Proposition~\ref{finite:procedure-soundness}
therefore applies. Table~\ref{finite:results} records output $1$ or $2$
for every enumerated pattern, each of which contradicts the supposed
realization. This proves the proposition.
\end{proof}

\begin{proof}[Proof of Theorem~\ref{thm:main}]
First suppose $m\ge n\ge3$ and $(m,n)\ne(3,3)$.
Then $m\ge4$, so $k=m-1\ge3$, and $2\le n-1\le k$.
If the upper bound failed, choose the least failed independence
parameter as in Section~\ref{sec:reduction}. The resulting graph $G$
has no $C_{k+1}$, satisfies $\alpha(G)\le k$, and has the expansion
property of Lemma~\ref{red:expansion}.
Theorem~\ref{clq:bound} gives
$\max\{3,\lfloor k/2\rfloor\}\le t=\omega(G)\le k$.
Proposition~\ref{small:cases} rules out $k=3,4$.
For $k\ge5$, Proposition~\ref{terminal:contradiction} rules out $t\ge9$.
The remaining values satisfy~\eqref{finite:domain}, so
Proposition~\ref{finite:verified} rules them out as well.
This proves the upper bound.

For the lower bound, colour the edges within each of $n-1$ disjoint
sets of size $m-1$ red and all edges between sets blue. A red connected
graph lies in a single set and hence cannot be a $C_m$. A blue clique
uses at most one vertex from each set and hence has order at most
$n-1$. This colouring has $(m-1)(n-1)$ vertices, proving the matching
lower bound.

Finally, colour the edges of a $5$-cycle red and its complementary
$5$-cycle blue. This gives $R(C_3,K_3)\ge6$. In a two-colouring of
$K_6$, one vertex has at least three neighbours of one colour, say red.
Any red edge among those neighbours completes a red triangle; if there
is none, the three neighbours form a blue triangle. Thus
$R(C_3,K_3)=6$.
\end{proof}
